\scanpage{59}
en prenant encore, comme dans 5
\[
(14)\qquad\rho_\infty=\tau_1\tau_0,\quad\rho_1=\tau_0\tau_\infty,\quad\rho_0=\tau_\infty\tau_1.
\]
On a donc
\[
(15)\qquad G_2\simeq\unsure{\Pi_{03}^\tau|\widetilde\Pi_{03}^\tau}/\langle\rho_1^2\rangle,
\quad G_2^\circ\simeq\Pi_{03}/\langle\rho_1^2\rangle
\]
\deleted{Cette fois} \inserted{Ici} \deleted{\unsure{G}} $\Pi_{03}^\tau$ apparaît comme le groupe ``classifiant'' pour la structure de ``\uline{carte tri-angulaire pondérée}'' (par $\{0,1,\infty\}$), i.e. une carte $C=(X,K)$ au sens de \unsure{p.g 1}, avec \deleted{l'ensemble} une application
\diagram{\begin{tikzcd}[column sep=large]
(16)\quad\begin{array}{c}S\\[-.2em]\text{ens des}\\\text{sommets}\\\text{(ordinaires)}\\\text{de la carte (\uline{pas} d'une}\\\text{subd. barycentrique!)}\end{array}\arrow[r]&\{0,1,\infty\}
\end{tikzcd}}
ayant la propriété que pour toute face $f$, l'application (16) induit une bijection de l'ens des sommets de $f$ avec $\{0,1,\infty\}$. L'ensemble des \uline{repères} (au sens de la théorie \inserted{tri-angulaire} pondérée, en cas de possibilité de confusion on dira ``repères pondérés'') de $X$ est par définition l'ens de ses faces, et $\unsure{\sigma_i|\tau_i}$ ($i\in\{0,1,\infty\}$) est le \deleted{\unsure{symé}} repère symétrique p.r. au \deleted{\unsure{repère le}} côté du repère opposé au \deleted{le} sommet de type $i$. L'\uline{ordre} des \uline{sommets} est pair on appelle \uline{ordre réduit} d'un sommet le quotient de l'ordre par \deleted{\unsure{2}} deux.

\mspage{58}
